% !TeX root = main.tex
\chapter{Prompts Used in the Zero-Shot Solution}
\label{appendix:zero-shot-prompts}

\section{System Prompts}

\begin{promptbox}
Java-to-Python:
You translate competitive-programming programs from token-spaced java into valid python. Return only python source code.

Python-to-Java:
You translate competitive-programming programs from token-spaced python into valid java. Return only java source code.
\end{promptbox}

\section{Base User Prompt Template}

\begin{promptbox}
Translate this token-spaced {src_lang} program into normal {tgt_lang} source code.

Important:
- Reconstruct the source program before translating it.
- NEW_LINE means a line break. INDENT and DEDENT describe Python-style block nesting in the source.
- Preserve the algorithm, constants, input/output behavior, and edge cases.
- Write complete runnable {tgt_lang} code.
- Do not copy {src_lang} libraries, methods, wrappers, or syntax into the answer.
- Output normal {tgt_lang} source code, not a token stream.
- Keep hardcoded values hardcoded. Only read input when the source program reads input.
- Mentally check the final code for valid {tgt_lang} syntax before returning it.

Target language: {tgt_lang}
{target_examples}

Return only the translated {tgt_lang} source code. Do not explain. Do not use Markdown.

Source tokenized {src_lang} program:
{code}
\end{promptbox}

\section{Python Target Prompt Addition}

\begin{promptbox}
Python output requirements:
- Output Python code only.
- Do not use Java syntax such as public, static, void, class Main, Scanner, String[] args, System.out.println, System.out.printf, charAt, toCharArray, or semicolons as statement endings.
- Use Python input(), sys.stdin, print(), sys.stdout.write(), lists, strings, math, and normal Python functions/classes where appropriate.
- Translate Java string methods to Python: s.charAt(i) becomes s[i], s.length() becomes len(s), toCharArray() becomes the same string or list(s).
- Translate Java Math methods to Python math functions or operators.
- If Java uses geometry/library objects, reimplement the math with Python functions, tuples, or complex numbers. Do not call methods like intersects(), distance(), p1, or p2 on tuples.
- Avoid shadowing Python builtins such as len, str, input, sum, max, min, list, dict, and set.
\end{promptbox}

\section{Java Target Prompt Addition}

\begin{promptbox}
Java output requirements:
- Output Java code only.
- Do not use Python syntax such as def, print(), range(), while True, import sys, list methods, string zfill, or colon-based blocks.
- Use Java Scanner/BufferedReader/StringTokenizer only when the source reads input. Keep hardcoded test values hardcoded.
- Use System.out.print/println/printf, arrays, ArrayList, StringBuilder, Math, braces, and semicolons where appropriate.
- Translate Python ** to Java Math.pow(...) or direct multiplication; never output * *.
- Translate Python zfill(n) to zero padding with String.format(...).replace(' ', '0') or equivalent Java code.
- Initialize boolean sieve arrays correctly, usually with Arrays.fill(prime, true), before marking composites.
- Methods called from main should be static unless they are called through an object.
\end{promptbox}
